\section{The secondary update}
\label{sec:secondary}

We now separate the pairs left by the two tests in
Section~\ref{main:section}. Their endpoints are seeds, grouped by their
broad and narrow bands, their padded $g$-cell, and their two old chart
identifiers. Proposition~\ref{pack:bound} bounds by $M_{\mathrm{pack}}$
the number of these groups meeting a narrow-band ball of radius $4R$.
This lets us isolate one group's clipping at an endpoint with positive
probability. In Section~\ref{pad:assembly}, the same bound will control
the number of formulas competing on a small ball.

Clipping raises the field using a proposal based on distance to admitted
seeds. In an ordinary chart it may also create an ordered chart, whose
auxiliary field starts at zero. For a surviving pair in an ordered chart,
either an endpoint is admitted in some orientation, allowing clipping,
or each endpoint has short paths on which the fixed ordered part rises
and falls. In the latter case Lemma~\ref{ord:two-side} places the
endpoints on opposite arcs, and the increment mode separates them using
the transverse field. Clipping and increment are alternatives at the
current scale; a newly created ordered chart is available for increments
at subsequent scales.

\subsection{Clipping near admitted seeds}
Fix the orientation of the old $B$ chart. In an ordinary chart its
comparison field is $p_*=s$; in an ordered chart it is one of $P$ and $-P$,
chosen fairly and shared by the entire old identifier. For each group
choose a representative $w_c$ and set $p_c=p_*(w_c)$. The weak diameter
bound of its $g$-cell, together with old-chart eligibility, gives
\begin{equation}\label{sec:variation}
 |p_*(w)-p_c|\le60g
 \quad\text{at every seed }w\text{ of the group}.
\end{equation}
All short paths used for this comparison lie in the old field band.

In an ordinary chart choose an ordered support $L$ with
\begin{equation}\label{sec:support}
 L\le p_*,\qquad L(w_c)=p_c,\qquad \Lip_{\mathrm{vert}}(L)\le1,
\end{equation}
using Lemma~\ref{ord:support} with the source of $s$ and target $w_c$.
For a constant chart use $L=0$. All seeds in an ordinary chart are
\emph{admitted}. In an ordered chart admit $w$ in the chosen orientation
if
\begin{equation}\label{sec:admission}
 p_*(v)\ge p_c-0.65R
 \quad\text{whenever }d_{J_2}(w,v)=R.
\end{equation}
The test is vacuous when that sphere is empty. Admission supplies the
boundary comparison needed by the shallow proposal below. Failure in
both orientations will instead supply paths on which the comparison
field rises and falls.
These decisions use no fresh clipping noises.

A group's seeds have $d_{J_1}$-diameter at most $10g$, but may lie in
several components of the narrow band. Give the group one independent
fair coin, and keep its representative $w_c$, value $p_c$, orientation,
and support $L$ (when present) common to all those components. For a kept
group, let $F$ be its admitted seeds in one narrow-band component,
provided this set is nonempty, and put
\[
 d_F(v)=\min_{w\in F}d_{J_2}(w,v).
\]
Apply Theorem~\ref{int:interpolation} to $0.9d_F$ on the entire
component, with error $E_B$, obtaining $Q_F$. Use the resulting proposal
only on $U_F=\{d_F\le R\}$. Thus broad-band smallness controls the
variation within a group, while narrow-band distance controls where
its proposals can act. The shallow branch supplies the gain at seeds.
In an ordinary chart the ordered-support branch protects a boundary where
the old field decreases too fast. The height shield protects the band ends.
Propose on this domain the minimum of
\begin{align}
 T_F(v)&=p_c+0.05R-Q_F(v),\label{sec:shallow}\\
 O_F(v)&=p_c+0.05R+4(L(v)-p_c)+\alpha_F,
             \quad\text{in an ordinary chart only},\label{sec:ordered-branch}\\
 H_F(v)&=\psi_F(h(v)).\label{sec:height-branch}
\end{align}
Here $\alpha_F$ is uniform in $[0,E_B]$ and independent of the fresh
interpolation noises defining $Q_F$. The height shield $\psi_F$
is $p_c-1+\alpha_h$ at the two ends of $J_2$, rises linearly to
$p_c+2+\alpha_h$ within distance $0.09a_2$ of each end, and is constant
on the intervening plateau; $\alpha_h$ is independently uniform in
$[0,E_B]$. Extend it constantly beyond the band. Conditional on the old
charts, bands, groups, orientations, admissions, and group coins, these
interpolation and branch noises are fresh, and the noise families of
different groups are independent. Replace $p_*$ by its maximum with
the proposals of all kept groups on their domains, and convert back to
the old field orientation.

\begin{lemma}[Validity of secondary clipping]\label{sec:validity}
The secondary update preserves the nested-set and chart invariants.
Its movement at each point, measured in symmetric difference, is bounded
by a universal constant in scale units.
\end{lemma}
\begin{proof}
First consider a point with $d_F(v)=R$. The shallow branch satisfies
\[
 T_F(v)\le p_c+0.05R-0.9R+E_B=p_c-0.849R.
\]
In an ordered chart a nearest admitted seed gives
$p_*(v)\ge p_c-0.65R$, leaving a margin of $0.199R$.
In an ordinary chart the same margin applies whenever that last
inequality holds. Otherwise $L\le p_*$ gives
\[
 O_F(v)-p_*(v)
 \le0.05R+3(p_*(v)-p_c)+E_B<-1.899R.
\]
Thus the proposal has a uniform margin below the old field on every
distance boundary. At the height ends the shield gives a margin as well:
each domain point lies within $R$ of a seed, which lies within $10g$ in
$d$ of $w_c$, and the old comparison has Lipschitz constant less than $5$
along these eligible paths. Its difference from $p_c$ is less than
$5(R+10g)<1/2$.

Every domain lies in its old chart, by seed eligibility. The comparison
orientation and label variables are common on overlapping domains.
Extend each proposal, after clamping against the old field, by the old
field off its domain. The boundary margins and the clipping argument in
Section~\ref{sets:section} show that this extension glues vertically and
respects switches. The finite maximum of these extensions has the same
properties. The spatial branches have Lipschitz constant at most $4$;
the shield has constant $3/(0.09a_2)$, which is less than $1/\kappa_B$.

A winning atom of $Q_F$ gives a constant plus or minus one point-source
distance and hence a fresh ordinary chart on its narrow-band component.
A winning $O_F$ branch, after conversion back to the old orientation,
becomes the ordered part of a new chart on the old band, with $u=0$.
It has slope at most $4$. Additive constants preserve its order; if
conversion negates the comparison field, reverse the two boundary
columns as in Definition~\ref{ord:definition}. This ordered part, its
band, and its arc assignment remain fixed under the new identifier;
subsequent increments change only its auxiliary field.
A height branch is rebased by
Lemma~\ref{sets:vertical}. Unchanged points retain their chart, and uses
of the same inherited atom in the same interpolation keep one identifier.

Finally $Q_F$ lies within $E_B$ of $0.9d_F$. Thus the shallow branch is
at most $p_c+0.051R$ and the old field is within $5(R+10g)$ of $p_c$
where the group can act. Taking a maximum over any number of groups does
not add these bounds. Formula~\eqref{sets:distance} gives a uniform
symmetric-difference movement bound.
\end{proof}

\begin{lemma}[Separation at admitted seeds]\label{sec:admitted-success}
If a surviving pair has a seed admitted for one comparison orientation,
then the no-change and secondary-clipping modes together give a probability
bounded below by a positive universal constant of separating that pair by
a positive universal amount.
\end{lemma}
\begin{proof}
At an admitted seed $w$, $|Q_F(w)|\le E_B$. Equations
\eqref{sec:variation}--\eqref{sec:support} imply that every applicable
spatial branch exceeds $p_*(w)+0.04R$: for the ordered branch use
$|L(w)-p_c|\le10g$. The height shield is on its interior plateau.
The other pair endpoint is outside this group's update domains, because
the group's seeds have unrestricted diameter at most $10g$ and the pair
distance is at least one.

At either endpoint all groups that could act have a seed within $R$ in
the narrow band. At most $2M_{\mathrm{pack}}$ groups can therefore act in total at the two endpoints.
Keeping the designated group and suppressing all other such groups has
probability at least $2^{-(2M_{\mathrm{pack}}+1)}$, a convenient lower bound. The desired ordered orientation,
when needed, has probability $1/2$.

Let $B_x,B_y$ denote the old system-$B$ sets and $B'_x,B'_y$ the updated sets.
On this event one endpoint is unchanged, while the movement of the other
has measure at least $0.04\kappa_BR$. If
$\mu(B_x\mathbin\triangle B_y)\ge0.01\kappa_BR$, the no-change mode
already gives the required separation. Otherwise the triangle inequality
gives $\mu(B'_x\mathbin\triangle B'_y)>0.03\kappa_BR$, and in
particular at least $0.02\kappa_BR$.
Each mode has probability $1/4$, so both alternatives have uniform
positive probability.
\end{proof}

\subsection{The increment in ordered charts}\label{sec:increment}
For an ordered identifier, the ordered part $s$, its band $I$, and its
arc assignment stay fixed while the auxiliary field $u$ changes and the
current region $\mathcal R$ may shrink. Its transverse field $t_s$ is
therefore fixed as well. We define the increments in original length
units, so these stored fields do not change when a step is analyzed in
units of its scale.

Define a continuous trapezoidal function $H:[0,\infty)\to[0,1]$ which vanishes
off $[0.005R,10]$, equals one on $[0.01R,8]$, and is linear on the two
remaining intervals. At scale $r$, in original units, add to $u$ the
function
\begin{equation}\label{sec:increment-formula}
 \pm t_s(v) H(|t_s(v)|/r)
 \min\{1,d(v,X\setminus\mathcal R)/r\}.
\end{equation}
The sign is fair and common to the old identifier. When the complement is
empty its distance factor is one. The formula is defined on the chart's
whole band, with changes to the sets effective only in its current region.
Old terms of $u$ are retained without modification.

\begin{lemma}[Increment bound]\label{sec:increment-bound}
If $Q>2000/R$, all increments of one ordered identifier together have
$d_I$-Lipschitz constant at most $K=100$. They preserve the nested sets,
introduce no chart interfaces, and move each point by at most a universal
constant times the scale.
\end{lemma}
\begin{proof}
The function $t_s$ is $4$-Lipschitz in $d_I$. The scalar function
$t\mapsto tH(|t|/r)$ has Lipschitz constant at most $6$ and magnitude
at most $10r$. The distance factor has Lipschitz constant at most $1/r$
for $d$, hence for $d_I$. Each increment has Lipschitz constant at most
$24+10=34$.

Its support is contained in
$0.005Rr<|t_s|<10r$. Since $t_s$ is fixed for this identifier, the
support intervals at distinct scheduled scales are disjoint when
$Q>2000/R$. In a comparison of two points at most two summands are
nonzero at either endpoint. Subtracting their zero values at the other
endpoint if necessary bounds the difference by $68d_I(v,w)$. Thus
$\Lip(u)\le68<K$ at all times.

The distance factor makes the movement vanish continuously at the region
boundary. Inside the region the new field has Lipschitz constant at most
$4+\lambda K<5<1/\kappa_B$. The vertical gluing argument applies, and
the region partition does not change. The amplitude bound is $10r$ before
the factor $\lambda$, proving the last assertion.
\end{proof}

\begin{lemma}[The remaining pairs]\label{sec:remaining}
Suppose a surviving pair has no admitted endpoint in either ordered
orientation. Then the increment mode separates it by a positive universal
amount with uniformly positive probability.
\end{lemma}
\begin{proof}
Normalize the current scale to one. Ordinary charts admit every seed,
so the old $B$ chart is ordered. Failure of each admission test supplies
a point on the corresponding $d_{J_2}$-sphere of radius $R$ and an
attaining path of length $R$. Failure for $P$ supplies a decrease of
$P$, and failure for $-P$ supplies an increase. By
\eqref{sec:variation}, each change has magnitude greater than
$0.65R-60g$. The correction $\lambda u$ changes by at most
$\lambda KR$ on either path. Thus the ordered part $s$ has a positive
and a negative change of magnitude greater than
\[
 0.65R-60g-\lambda KR>0.4R
\]
along these paths from each endpoint.

A shortest $x$--$y$ path has length at most $2$ and stays in the eligible
radius-$5$ ball, hence in the old band $I$. This proves
$d_I(x,y)=d(x,y)$ and bounds the change of $u$ by $2K$. The $B$ cutoff
therefore gives
\[
 |s(x)-s(y)|\le10^{-4}R+2\lambda K<0.001R.
\]
The four excursion paths lie in the common band $J_2$, whose width
satisfies Lemma~\ref{ord:two-side}. Applying that lemma gives opposite
signs of $t_s$ and $0.01R\le |t_s(x)|,|t_s(y)|\le8$.

These bounds put both endpoints on the plateau of the trapezoidal
factor. Eligibility also puts their radius-$5$ balls in the current
region, so the distance factor in \eqref{sec:increment-formula} equals
one at both endpoints. For one of the two increment signs,
\[
 |P_{\mathrm{new}}(x)-P_{\mathrm{new}}(y)|
   \ge \lambda|t_s(x)-t_s(y)|\ge0.02\lambda R.
\]
Formula~\eqref{sets:distance} proves separation at least
$0.02\kappa_B\lambda R$. The favorable sign has probability at least
$1/2$.
\end{proof}
